% explainer-source-hash: eb78dcea1059a3fad077ff13dc88bb7062fcb0ae354b0341bf97777c95f2d653
% explainer-source-commit: df2470cbd060525eba573923539bceac51b89a13
% explainer-generated: 2026-07-16
% Regenerate with /explain-all (see WIRING.md). Do not edit this stamp by hand:
% it is how the toolchain knows whether this document still matches the code.
\documentclass[11pt,a4paper]{article}
\input{preamble}

\title{\vspace{-1.5cm}\textbf{Police Shootings Analysis}\\[0.3em]
\large A Brief: what it does, what it gets right, and what is broken}
\author{Explainer for \file{police-shootings-analysis}}
\date{}

\begin{document}
\maketitle
\thispagestyle{empty}

\section*{In one paragraph}

\file{police-shootings-analysis} reads five fixed CSV files and turns them into charts, two
different ways. A single 392-line script, \file{analyse\_deaths.py}, writes 24 static chart
files and exits. A small Flask web app, \file{app.py}, serves three of the same views as an
interactive dashboard that filters by state, year, and armed status. It began as a guided
Udemy exercise about joining messy real-world data and comparing three plotting libraries.
The dashboard is later, better work built beside it. The two halves share one CSV file on
disk and essentially no code.

\begin{honest}{The subject, and the headline finding}
Every row of the main dataset is a person who was killed. This document treats that data
descriptively and states its limits rather than repeating chart titles on trust.

The single most important thing to know about this project: \textbf{three of its 24 charts
are titled ``High School Graduation Rate by State'' and actually plot the number of cities
per state}, on a y-axis running to 1750. One of them is committed to the repository as a PNG.
The cause is \code{.count()} where \code{.mean()} was meant. The fix is one word. It is
covered in the ``What is broken'' section, and in full in the Deep Dive.
\end{honest}

\section{What the data can and cannot say}

\begin{table}[H]
\centering
\small
\begin{tabular}{@{}p{6.5cm}p{8.5cm}@{}}
\toprule
\textbf{What it is} & \textbf{What that means} \\
\midrule
A fixed extract of The Washington Post's fatal-police-shootings database: 2,535 records,
2015 through part of 2017.
& Not live, not current, and the code has no way to fetch an update. A chart titled
``Killings by Year'' stops in 2017 and cannot know that. \\
\addlinespace
Four city-level Census snapshots from around 2015 (poverty, high school completion, race
share, income).
& Static too. And the fatality data is never actually joined to any of them, despite the
README's framing. \\
\addlinespace
Counts of \emph{records}. & Not counts of events. Anything that shaped what got reported is
invisible here. \\
\addlinespace
Raw totals with no population denominator. & The map and the ``top 10 cities'' chart
substantially rank by population. California leads at 424 records and is the most populous
state. A code comment asks ``which cities are the most dangerous''; the chart cannot answer
that. \\
\addlinespace
Correlations. & Never causation. The strongest relationship in the project, city poverty
against city high school completion, is about $-0.50$, verified. That is an association and
nothing more. \\
\addlinespace
Some very small cells. & The race-by-top-10-city chart's bars hold between 10 and 21 records
each, verified. Nothing stable survives at that size. \\
\bottomrule
\end{tabular}
\end{table}

\section{Architecture}

\begin{figure}[H]
\centering
\resizebox{\textwidth}{!}{
\begin{tikzpicture}[node distance=8mm]
  \node[store] (deaths) {Deaths\_by\_Police\_US.csv\\2,535 rows};
  \node[store, below=16mm of deaths] (census) {3 Census CSVs used\\+ 1 loaded and ignored};

  \node[box, right=32mm of deaths, minimum width=26mm] (datapy) {data.py};
  \node[box, right=24mm of datapy, minimum width=24mm] (apppy) {app.py\\\scriptsize Flask};
  \node[box, right=24mm of apppy, minimum width=26mm] (html) {index.html\\\scriptsize Plotly.react};

  \node[box, right=32mm of census, minimum width=32mm] (script) {analyse\_deaths.py};
  \node[store, right=42mm of script] (out) {output/\\18 PNG + 6 HTML};

  \draw[flow] (deaths) -- (datapy);
  \draw[flow] (datapy) -- node[lbl,above,align=center]{clean once\\at startup} (apppy);
  \draw[flow] (apppy) -- node[lbl,above,align=center]{figure JSON} (html);
  \draw[flow] (census) -- (script);
  \draw[flow] (script) -- node[lbl,above,align=center]{savefig /\\write\_html} (out);
  \draw[flow] (deaths.south west) to[out=205,in=155] node[lbl,left,align=center]{re-read and\\re-cleaned\\independently} (script.west);

  \begin{scope}[on background layer]
    \node[fit=(datapy)(apppy)(html), draw=accent, dashed, rounded corners,
          inner sep=5mm, label={[lbl]above:the dashboard: careful code}] {};
    \node[fit=(script)(out), draw=inkblue, dashed, rounded corners,
          inner sep=4mm, label={[lbl]below:the original exercise: unreviewed}] {};
  \end{scope}
\end{tikzpicture}}
\caption{Two entry points sharing a file, not a code path. The curved arrow is the whole
problem: \file{analyse\_deaths.py} does not import \file{data.py}, so the two halves keep
duplicate copies of the same cleaning logic and can drift apart.}
\end{figure}

\begin{keyidea}{The shape to remember}
\file{data.py} is described in its own docstring as the single home of the cleaning logic. It
is the single home of that logic \emph{for the dashboard}. \file{analyse\_deaths.py} still
has its own inline copy and imports nothing. The abstraction was extracted for the new code;
the old code was left pointing at a duplicate.
\end{keyidea}

\subsection{The static script}

Straight-line code, no classes, two small functions. Four decisions define it:

\begin{enumerate}
  \item \code{matplotlib.use('Agg')} on line 2, before \code{pyplot} is imported. This is the
  single most valuable change made to the original notebook exercise: it makes the script
  run anywhere, with no display, writing durable files instead of opening windows that never
  appear.
  \item A global counter, \code{next\_chart\_path()}, names every output \code{chart\_NN}. The
  numbering is pure execution order, so the PNGs are only meaningful alongside the exact
  script version that made them, and nothing enforces that pairing.
  \item \code{fillna(0)} on all five dataframes, before any numeric coercion. This one is a
  bug and is covered below.
  \item The same chart drawn in matplotlib, seaborn, and plotly, which is the actual point of
  the exercise.
\end{enumerate}

Charts 01--12 come from the Census data. Charts 13--24 come from the fatality data. The two
groups never meet.

\subsection{The dashboard}

\file{data.py} loads and cleans once and exposes \code{apply\_filters()}. \file{app.py} holds
the cleaned dataframe in memory at import time, and on each request filters it and rebuilds
three Plotly figures (a choropleth, a year bar chart, a top-10-cities bar chart) plus a
four-value stat row. \file{templates/index.html} server-renders the first payload, then calls
\code{GET /api/charts} on every filter change and updates the charts in place with
\code{Plotly.react}, so hover and zoom survive and nothing is a static image. Clicking a
state on the map sets the state dropdown and re-fetches, so the map and the dropdown cannot
disagree.

\section{Verified by running it}

Everything here came from a clean virtual environment built from \file{requirements.txt}, not
from the README.

\begin{table}[H]
\centering
\small
\begin{tabular}{@{}p{6.2cm}p{5.4cm}p{3.4cm}@{}}
\toprule
\textbf{Claim} & \textbf{Result} & \textbf{Verdict} \\
\midrule
2,535 records; 2,428 M / 107 F & as stated & confirmed \\
633 of 2,535 with signs of mental illness & 25.0\% & confirmed \\
Top armed: gun 1,398, knife 373, vehicle 177 & as stated & confirmed \\
Top states: CA 424, TX 225, FL 154 & as stated & confirmed \\
24 charts produced (18 PNG, 6 HTML, 31 MB) & as stated & confirmed \\
Dashboard, including the zero-row filter case & renders, shows n/a & confirmed \\
\midrule
Killings per year & \textbf{991 / 963 / 581}, not 991 / 962 / 582 & \textbf{README wrong} \\
``regenerates in under a minute'' & \textbf{6 min 54 s} & \textbf{README wrong} \\
\bottomrule
\end{tabular}
\end{table}

\section{What is broken}

\begin{honest}{1. Three charts plot the wrong quantity (charts 04, 05, 06)}
\code{groupby('Geographic Area')['percent\_completed\_hs'].count()} counts rows per state. It
never looks at the column's values. \code{.mean()} is what a rate needs.

So three charts titled ``High School Graduation Rate by State'' plot how many cities each
state has. Rendered and read: the y-axis reaches 1750, the tallest bars are PA 1,761, TX
1,710, CA 1,501, and the shortest is DC at 1, because DC contributes one city row. Chart 04
is committed to the repository.

The tell was available: chart 07, in the same script, plots the same quantity correctly via
seaborn's implicit mean, topping out near 100. Two charts in one document disagree by a
factor of 17. Nothing noticed, because the project has no tests at all.
\end{honest}

\begin{honest}{2. The income dataset is loaded and never used, and nothing is cross-referenced}
The README's first sentence says the fatality data is ``cross-referenced against four US
Census datasets $\ldots$ to look for socioeconomic patterns''.

\code{df\_hh\_income} appears twice in the whole project: line 37 reads the 709 KB file, line
46 calls \code{fillna(0)} on it. Nothing reads it again. Verified by grep.

The larger version, which an interviewer will reach: \textbf{the fatality data is never
joined to any Census data at all}. The Census charts and the fatality charts are two separate
analyses sharing a script. The cross-reference is an aspiration, not a feature.
\end{honest}

\begin{honest}{3. \code{fillna(0)} runs before the coercion meant to catch bad values}
Because missing cells become the number 0 first, the \code{errors='coerce'} that was supposed
to turn them into NaN never sees them. Verified consequences: \textbf{77 records with no
recorded age are plotted as zero-year-olds} in the age histogram and KDE (chart 17), pulling
the density curve left with no visual cue that the points are fabricated; and \textbf{195
records with no recorded race become a bar literally labelled \code{0}} in chart 18.

The inline comment above those lines reads ``Replace 0 with NaN''. The code does the exact
opposite, and the two operations have opposite consequences.

\file{data.py} inherited the same inverted order, so both halves share it. In the dashboard's
favour: \file{app.py} and \file{data.py} both explicitly strip the \code{0} values at the
point of use, so the author did notice the symptom there and patched it locally rather than
at the source. The script never got that treatment.
\end{honest}

\begin{honest}{4. Chart 12 is not stacked, and the README says it is}
The loop draws five bar series at \code{alpha=0.7} and never passes \code{bottom=}, so all
five start at zero and overplot each other. Rendered and confirmed: the White bar covers
almost everything drawn after it. A reader who expects a stacked chart reads a bar at 92 as a
cumulative total when it means ``the White share is 92''.

The underlying arithmetic is also an unweighted mean across cities, so a village counts as
much as a city of 800,000. Defensible, but never acknowledged.
\end{honest}

\begin{honest}{5. \code{scatter=\{'alpha': ...\}} does nothing (charts 10, 11)}
Verified against the installed seaborn 0.13.2 signature: \code{regplot} and \code{lmplot}
take \code{scatter} as a \textbf{boolean}, and \code{scatter\_kws} as the styling dict. A
non-empty dict is truthy, so the alpha is silently discarded and no error is raised. The
author knew the right pattern: \code{line\_kws=\{'color': 'red'\}} sits right beside it and
works correctly. So two scatter plots of 29,329 points render at full opacity when
transparency was intended.
\end{honest}

\subsection{Smaller ones}

\begin{itemize}
  \item \textbf{\code{convert\_country\_code} is a guarded no-op}, kept alive purely to
  support chart 23, a world-scope map plotting US state postal codes onto countries. Chart 23
  is not meaningful and \code{pycountry} exists in \file{requirements.txt} only for it. (That
  some state codes collide with real ISO country codes and are therefore silently rewritten
  follows from \code{pycountry}'s contract but was \textbf{not measured} on this dataset:
  treat it as inferred.)
  \item \textbf{Chart 01 does not aggregate.} matplotlib is handed 29,329 city rows and draws
  them all onto 50 x-positions. Chart 02, the seaborn version of the same title, averages per
  state. Two charts, same title, different quantities.
  \item \textbf{\code{GET /} ignores its query string}, so a filtered view has no URL and
  cannot be bookmarked or shared. \code{/api/charts} already parses exactly those parameters
  correctly.
  \item \textbf{The Plotly CDN} is an unpinned single point of failure with no integrity
  hash, so a local dashboard over local data needs the internet to draw anything.
  \item \textbf{\code{fetchAndRender} has no \code{catch}}, so a failed request leaves stale
  charts on screen that no longer match the dropdowns.
  \item \textbf{No tests and no schema validation.} This is the root cause of items 1 through 5.
\end{itemize}

\section{What it gets right}

\begin{itemize}
  \item The headless \code{Agg} backend plus file output is what makes the script
  reproducible at all. It was a real fix to a real problem.
  \item \code{base\_path} from \code{sys.argv[0]}: the script finds its data from any working
  directory.
  \item Exact version pins on every dependency, including \code{scipy}, which the project
  \emph{never imports} and cannot run without: both KDE charts delegate to
  \code{scipy.stats.gaussian\_kde} lazily at draw time. Declaring a behavioural dependency
  the import graph does not show is exactly right.
  \item \code{format='mixed'} on the date column, which genuinely mixes \code{M/D/YYYY} and
  \code{DD/MM/YY}. Verified to parse all 2,535 rows with zero failures.
  \item \code{fig\_to\_json} routes Plotly figures through Plotly's own encoder, sidestepping
  a real Flask/numpy serialization failure before it could happen at request time.
  \item \code{build\_stats} handles the empty selection. Tested: it returns nulls and the page
  shows ``n/a'' instead of raising \code{IndexError}.
  \item The \code{\# noqa: E712} on \code{== True} is the author correctly overruling the
  linter: on a pandas Series that is a vectorised comparison and \code{is True} would be wrong.
  \item Map clicks route \emph{through} the dropdown rather than around it, so two controls
  over one piece of state cannot disagree.
  \item The README's ``Data and framing notes'' and the dashboard footer both state the
  causation and sample-size limits unprompted. On this subject that matters, and most
  exercise write-ups skip it.
\end{itemize}

\section{The fix list}

\begin{enumerate}
  \item \code{.count()} to \code{.mean()} at line 89. One word, fixes three false charts.
  \item Move \code{fillna(0)} after the coercions, or delete it. Fixes charts 17 and 18, in
  both entry points.
  \item Use \code{df\_hh\_income} or delete it, and rewrite the README's ``four Census
  datasets $\ldots$ cross-referenced'' to describe what the code does.
  \item Pass \code{bottom=} in the chart 12 loop, or stop calling it stacked.
  \item \code{scatter=} to \code{scatter\_kws=} in charts 10 and 11.
  \item Delete \code{convert\_country\_code}, chart 23, and \code{pycountry}.
  \item Correct the README's year counts and its runtime claim.
  \item Give each chart a function and a name; add one test that asserts a rate chart's
  maximum is under 100.
\end{enumerate}

\begin{keyidea}{The defence, in one line}
The dashboard half is careful, well-factored code with real thought in it. The static-script
half is an exercise that was made to run headlessly and then never looked at again, and it
contains three charts that are confidently, visibly false. The gap between the two is the
whole story, and the reason for the gap is that neither half has a single test.
\end{keyidea}

\section*{Glossary}

\begin{description}[leftmargin=!,labelwidth=3.0cm,style=nextline]
  \item[CSV] Plain-text table, commas between columns. Carries no type information, which is
  why a column of numbers arrives as text and has to be coerced.
  \item[Character encoding] The byte-to-character mapping. All five files here are
  Windows-1252, not UTF-8, and three of them use bare carriage returns as line endings.
  \item[DataFrame] pandas' in-memory table. The unit of work for everything in this project.
  \item[NaN / NaT] pandas' ``missing number'' and ``missing date''. Deliberately contagious,
  so that missingness forces a decision instead of quietly becoming a plausible number.
  \item[Coercion] \code{pd.to\_numeric(col, errors='coerce')}: convert to a number, turn
  anything unconvertible into NaN rather than crashing.
  \item[groupby] Split a table by a key, compute per group, reassemble one row per group.
  \item[Backend (matplotlib)] What matplotlib draws onto. \code{Agg} draws into memory
  instead of a window, which is what makes the script runnable with no display.
  \item[KDE] Kernel Density Estimate: a smooth curve estimating a distribution's shape. A
  smoothed guess, not a measurement, and it will smooth fabricated data into what looks like
  a real feature.
  \item[Choropleth] A map that colours regions by a value. Colours by count here, not rate.
  \item[Correlation] A number in $[-1, +1]$ for how tightly two quantities move together. Says
  nothing about why.
  \item[Flask] A minimal Python web framework. Maps URLs to functions.
  \item[Jinja2 / \code{| safe}] Flask's template language, and the filter that disables its
  automatic HTML escaping. Necessary for injecting pre-built JSON; safe here only because the
  data is trusted and static.
  \item[Serialization] Turning an in-memory object into bytes for the network. JSON knows few
  types, so numpy integers need an explicit rule.
  \item[Transitive dependency] A package installed because something else needed it. Relying
  on one without declaring it, as this project would with \code{scipy}, breaks the day the
  intermediate drops it.
  \item[CDN] A third-party host serving a library over the internet. Convenient, and a runtime
  dependency on someone else's uptime.
\end{description}

\end{document}
